\documentclass[11pt,a4paper]{article}
\usepackage[T1]{fontenc}
\usepackage{amsmath,amssymb,geometry,booktabs}
\geometry{margin=27mm}
\title{EPPA--III: remaining Lean obligations for the direct clique-cover proof}
\author{Independent validation supplement; original manuscript frozen}
\date{11 October 2026}
\begin{document}
\maketitle

\paragraph{Scope.}
The source of \emph{Graphs with Small EPPA Witnesses} remains
unchanged. This note identifies the exact residual graph-theoretic
obligations after the orbit-count repairs to Claim~3.9.
It distinguishes verified theorems from currently pending PRs.

\paragraph{Current formal boundary.}
The following statements have been axiom-audited and merged:
the EPPA neighbourhood-profile orbit upper bound;
the unrestricted \(S_t\wr S_s\) orbit lower bound for a fixed
coordinate size \(k\); the inclusion of this wreath orbit in the
actual EPPA profile orbit; the extreme-degree conclusion
\(k=1\) or \(k=t-1\); saturation injectivity of the
size-\(k\) coordinate projection; and the repeated-partial-size
obstruction forcing \(s=t=2\).
These are in the source modules
\texttt{ProfileOrbitBound},
\texttt{WreathOrbitCount},
\texttt{Claim39OrbitBridge},
\texttt{ProfileSaturation}, and
\texttt{RepeatedPartialSize}.

The independent-inner-movement theorem (PR~\#38),
the unique-size EPPA consequence (PR~\#39),
and the two-partial-coordinate exception (PR~\#40)
are separate stacked formalizations,
\emph{not counted as checked until their CI and axiom
audits succeed}. Together they aim to prove:
\[
 \begin{gathered}
 s,t\ge2,\quad H\supseteq_{\mathrm{ind}}sK_t,\quad
 |H|=2st,\quad H\text{ EPPA for }sK_t;\\
 \exists w\in H\setminus G,\quad
 \exists i\ne j,\quad
 0<|N(w)\cap C_i|,\ |N(w)\cap C_j|<t\\
 \Longrightarrow\quad s=t=2.
 \end{gathered}
\]
Neither transitivity nor an outside clique partition is needed
for this implication.

\paragraph{Next obligation A: equal-edge-count identity.}
For a vertex-transitive graph \(H\) on \(2st\) vertices, with
an induced copy \(G=sK_t\) and complementary part \(B\) of
the same order, regularity gives
\[
 |E(H[B])|=|E(H[G])|=s\binom t2.
 \tag{A}
\]
This is a basic finite degree-sum and cross-edge identity;
the graph-level Lean statement is still to be written.

\paragraph{Next obligation B: one partial coordinate.}
If a profile has exactly one partial coordinate of size \(k\),
the group action implies
\[
 |W\cdot P|=s\binom tk\binom{s-1}{r},
\]
where \(r\) counts the remaining full coordinates.
Together with the verified upper bound \(st\), one gets
\(k\in\{1,t-1\}\), \(r=0\) or \(s-1\), and
\(|W\cdot P|=st\). Hence all outside profiles are distinct
and exhaust the complementary half.
The resulting action on \(B\) is the standard imprimitive
action on \(s\) blocks of \(t\) vertices.
Its unordered pair orbits are the within-block
and across-block pairs. Invariance of \(E(H[B])\),
followed by (A), gives \(H[B]=sK_t\);
the cross-edges are the required matching/co-matching
diagonal with uniform off-diagonal entries.
The exact action identification and pair-orbit argument,
not merely the numerical restriction, are outstanding.

\paragraph{Next obligation C: two partial coordinates.}
Under the formally targeted \(s=t=2\) exception,
all four profiles are the singleton choices in both
selected \(K_2\)'s, indexed by the corners of a square.
The unordered pair-orbits of the induced wreath action
have sizes four (sides) and two (opposite corners).
By (A), \(B\) has exactly two edges, forcing the opposite
matching. Its cross-connection matrix then normalises to
\[
 \begin{pmatrix}
 M&M\\M&\overline M
 \end{pmatrix},
\]
the Wagner graph. This last eight-vertex orbit and
edge-count implication is a small formalization target.
The earlier Lean theorem excluding the cube remains available.

\paragraph{Next obligation D: no partial coordinates.}
All vertices within each \(C_i\) have equal closed
neighbourhoods, hence belong to one true-twin class.
Vertex-transitivity makes all twin classes have
the same order \(a\ge t\). The quotient \(Q\)
on twin classes has \(M=2st/a\le2s\) vertices
and contains \(s\) independent vertices.
For non-homogeneous \(H\), the quotient is
a positive-degree vertex-transitive graph.
Double-counting edges out of an independent set
in a positive-degree regular graph gives
\(\alpha(Q)\le M/2\).
Thus \(M=2s\), \(a=t\).
The \(s\) complementary twin classes already contain
\(s\binom t2\) internal edges; (A) excludes additional
edges, giving \(B=sK_t\).
Finally, EPPA and the orbit count on the
\(s\) complete/empty clique coordinates imply
that all outside-clique neighbourhoods are uniform,
or are exactly the singleton/co-singleton index orbit.
These are precisely the uniform cross-block patterns
of Definition~1.6. This quotient step is paper-level,
not currently Lean-verified.

\paragraph{Next obligation E: low parameters and homogeneous hosts.}
The cases \(s=1\) or \(t=1\) admit direct small-profile
proofs, and homogeneous \(H\) is handled separately
using the finite homogeneous-graph classification.
These are included at paper level in
\texttt{direct-clique-cover-proof.tex}; they are not
included in the statement of the two-partial-coordinate
Lean obstruction, which assumes \(s,t\ge2\).

\paragraph{Recommended validation order.}
Complete PRs~\#38--\#40 with green Lean CI and
permitted-axiom audit; prove (A), then (B) and (C),
and finally the regular twin-quotient step (D).
This establishes the direct proof of Lemma~3.7
without reusing the unverified \(D_j\)-regularity
step in Claim~3.8 or the simultaneous matching
normalisation in Claim~3.10.
Only then mark the graph theorem as fully verified.
The author-facing minimal Wagner correction is unchanged.

\end{document}
